% from ega4-18-fr.tex
\oldpage[IV-4]{167}

Soit alors $A'={}^{hs}\!A$, qui est un anneau local et un $A$-module fidèlement plat (18.8.8), l'homomorphisme $A\to A'$ étant local. Posons $Y'=\Spec(A')$, $X'=X\times_Y Y'$, $f'=f_{(Y')}:X'\to Y'$, de sorte que $f'$ est étale et séparé; en outre, comme le morphisme $Y'\to Y$ est plat, le raisonnement de (6.15.4.1) montre que $f'$ est aussi birationnel. Si $y'$ est le point fermé de $Y'$, il suffira, en vertu de (2.6.1, (v)), de prouver que $f'^{-1}(y')$ est radiciel sur $k(y')$; puisque $f'$ est étale, il suffira de prouver que $f'^{-1}(y')$ contient au plus un point (17.6.1, c')).

Supposons donc $Y=\Spec(A)$, $A$ étant strictement local, $f$ séparé, étale et birationnel, et montrons que $f^{-1}(y)$ ne peut contenir plus d'un point. En effet, s'il existait dans $f^{-1}(y)$ deux points distincts $x_1$, $x_2$, il existerait deux $Y$-sections $u_1$, $u_2$ de $f$ telles que $u_1(y)=x_1$ et $u_2(y)=x_2$ (18.8.1); mais puisque $f$ est étale et séparé, $u_1$ et $u_2$ seraient des isomorphismes de $Y$ sur deux composantes connexes (ouvertes) de $X$ (17.9.4) et il y aurait donc deux points maximaux au moins de $X$ au-dessus de tout point maximal de $Y$, ce qui contredit l'hypothèse que $f$ est birationnel.

\begin{proposition}[18.10.19]
\label{IV.18.10.19}
Soient $Y$ un préschéma réduit dont l'ensemble des composantes irréductibles est localement fini, $f:X\to Y$ un morphisme séparé, localement de type fini et formellement non ramifié, $g$ une $Y$-section rationnelle de $f$ (c'est-à-dire une $Y$-application rationnelle de $Y$ dans $X$), $U$ le domaine de définition de $g$ (I, 7.2.1), et soit $Z$ l'adhérence de $g(U)$ dans $X$. Alors, pour tout point $y\in Y-U$ tel que $Y$ soit géométriquement unibranche au point $y$, on a $Z_y=Z\cap f^{-1}(y)=\varnothing$. En particulier, si $Y$ est géométriquement unibranche, $g(U)$ est une partie à la fois ouverte et fermée de $X$. Si, pour toute partie fermée irréductible $T$ de $X$, $f(T)$ est fermé dans $Y$, $g$ est définie en tout point géométriquement unibranche de $Y$.
\end{proposition}

Puisque $Y$ est réduit et $f$ séparé, il résulte de (I, 7.2.2) qu'il existe une $U$-section $u$ de $f^{-1}(U)$ appartenant à la classe $g$; d'ailleurs (17.4.1.2), puisque $f$ est formellement non ramifié, $u$ est un isomorphisme de $U$ sur le préschéma induit sur l'ouvert $u(U)$ de $X$. Si on note encore $Z$ le sous-préschéma réduit de $X$ ayant $Z$ pour espace sous-jacent, $Z|_U=u(U)$, étant réduit, est induit aussi par $Z$, donc la restriction $f'=f|_Z$ de $f$ est un morphisme birationnel de $Z$ dans $Y$, d'ailleurs formellement non ramifié, puisque $f$ l'est (17.1.3). La question étant locale sur $Y$, puisque $Y$ est réduit et géométriquement unibranche (donc intègre) au point $y$, on peut se borner (en remplaçant $Y$ par un voisinage ouvert de $y$) au cas où $Y$ est intègre; alors $u(U)$ est irréductible, donc aussi $Z$, et par suite $Z$ est aussi intègre. Alors (I, 8.2.7), puisque $f'$ est dominant, l'homomorphisme $\sh{O}_{Y,f(x)}\to\sh{O}_{Z,x}$ est injectif pour tout $x\in Z_y$, et on conclut de (18.10.2, (iii)) que $f'$ est un morphisme étale; étant séparé et birationnel, c'est un isomorphisme local au point $x$ en vertu de (18.10.18); mais cela entraînerait que la $U$-section $u$ serait prolongeable à un ouvert distinct de $U$, contrairement à la définition de ce dernier. On a donc nécessairement $Z_y=\varnothing$, ce qui établit la première assertion; la seconde en est une conséquence évidente. En outre, sous les hypothèses de la dernière assertion de l'énoncé, comme $Z$ est fermé et irréductible, $f(Z)$ est fermé dans $Y$, donc égal à $Y$, c'est-à-dire $Z_y\ne\varnothing$ pour tout $y\in Y$, ce qui achève la démonstration.

\begin{remark}[18.10.20]
\EGAReaderTarget{ega4.c04.rem.18-10-20}{IV.18.10.20}
Soit $Y$ un préschéma localement noethérien. On dit qu'un morphisme $f:X\to Y$ est essentiellement propre s'il est localement de type fini et si, pour tout $Y$-schéma de la forme $Y'=\Spec(A)$ où $A$ est un anneau de valuation discrète, l'application canonique (II, 7.3.2.2) relative au $Y'$-préschéma $X\times_Y Y'$ est bijective.

